%----------------------------------------------------------------------------------------
% MACHTEN EN WORTELS. Alle opgaven zijn eigen werk.
%
% Opbouw: Start, leerstof (theorie, werkwijze of uitleg, voorbeeld, korte oefening), en op een
% nieuwe pagina de Opdrachten (door elkaar, stijgende moeilijkheid), Koppel en Test jezelf.
% Labels van dit hoofdstuk beginnen met h4-.
%----------------------------------------------------------------------------------------

\mychapter{4}{Machten en wortels}

\lkr{Leerplan B+S'': het tweede deel van LPD 42 (machtsverheffing en $n$-de machtsworteltrekking in goniometrische vorm), met de formule van De Moivre en de wenk over de meetkundige betekenis: de $n$-de machtswortels vormen een regelmatige $n$-hoek, opeenvolgende machten liggen op een spiraal. Het bewijs van De Moivre is een herhaling van de productregel; een bewijs met volledige inductie kan in de cursus Logica en bewijzen.}

Een hoge macht als $(1 + i\sqrt3)^6$ uitrekenen met het binomium is veel werk. In goniometrische vorm is het bijna niets: een macht is een product van gelijke factoren, dus je vermenigvuldigt de modulus een paar keer met zichzelf en telt het argument een paar keer op. Omgekeerd kun je zo ook wortels trekken. Dan blijkt dat elk complex getal $z \neq 0$ precies drie derdemachtswortels heeft, vier vierdemachtswortels, en zo verder, en dat ze altijd mooi regelmatig in het vlak liggen.

\begin{oefeningen}[Start]
\oef{h4-start-1}

\oef{h4-verken}
\end{oefeningen}

\lkr{Laat oefening \ref{oef:h4-verken} maken vóór de uitleg. In a en b ontdekken leerlingen de formule van De Moivre. In c raden ze een derdemachtswortel van $i$ door terug te rekenen; dat er nog een tweede (en een derde) is, is de verrassing waarmee sectie 2 begint.}

%========================================================================================
\section{De formule van De Moivre}

\subsection{Machten in goniometrische vorm}

Neem je in de productregel twee keer hetzelfde getal $z = r(\cos\theta + i\sin\theta)$, dan vind je $z^2 = r^2(\cos 2\theta + i\sin 2\theta)$. Vermenigvuldig je nog eens met $z$, dan komt er weer een factor $r$ bij de modulus en een $\theta$ bij het argument. Zo gaat het voor elke exponent.

\begin{theorie}
\begin{stelling}{Formule van De Moivre}
Voor $z = r(\cos\theta + i\sin\theta)$ en elk natuurlijk getal $n$ is

\vspace{1mm}
\qquad $\formule{z^n = r^n(\cos n\theta + i\sin n\theta)}$
\end{stelling}

\begin{bewijs}
Voor $n = 0$ staat er links en rechts 1, want $\cos 0 + i\sin 0 = 1$. Geldt de formule voor een exponent $n$, dan geeft de productregel
\bewijsregel{z^{n+1} = z^n \cdot z = r^n \cdot r\bigl(\cos(n\theta + \theta) + i\sin(n\theta + \theta)\bigr) = r^{n+1}\bigl(\cos(n+1)\theta + i\sin(n+1)\theta\bigr) ,}
dus geldt ze ook voor $n + 1$. Zo volgt ze uit $n = 0$ voor $n = 1$, daaruit voor $n = 2$, en zo verder voor elk natuurlijk getal.
\end{bewijs}
\end{theorie}

De formule draagt de naam van Abraham de Moivre, een Franse wiskundige die in Londen leefde. Ze zegt: de modulus van $z^n$ is $r^n$, en $n\theta$ is een argument van $z^n$. Ze geldt ook voor negatieve gehele exponenten, zoals je in opdracht \ref{oef:h4-o12} nagaat.

\voorbeeld{}{Bereken $(1 + i\sqrt3)^6$.\label{vb:h4-moivre}}{%
\uitleg{We zetten het grondtal in goniometrische vorm, en passen De Moivre toe: de modulus tot de zesde, het argument maal zes.}
\begin{stappen}
1 + i\sqrt3 = 2\Bigl(\cos\frac{\pi}{3} + i\sin\frac{\pi}{3}\Bigr) & $r = \sqrt{1 + 3} = 2$ en $\theta = \frac{\pi}{3}$. \\[3ex]
(1 + i\sqrt3)^6 = 2^6\bigl(\cos 2\pi + i\sin 2\pi\bigr) & $6 \cdot \frac{\pi}{3} = 2\pi$. \\[1ex]
= 64 & $\cos 2\pi = 1$ en $\sin 2\pi = 0$.
\end{stappen}}

\voorbeeld{}{Bereken $(1 - i)^{10}$.\label{vb:h4-moivre2}}{%
\uitleg{Ook hier eerst de goniometrische vorm. Het argument maal tien valt buiten $]-\pi, \pi]$: we brengen het terug met een veelvoud van $2\pi$.}
\begin{stappen}
1 - i = \sqrt2\Bigl(\cos\Bigl(-\frac{\pi}{4}\Bigr) + i\sin\Bigl(-\frac{\pi}{4}\Bigr)\Bigr) & $r = \sqrt2$, vierde kwadrant. \\[3ex]
(1 - i)^{10} = \bigl(\sqrt2\bigr)^{10}\Bigl(\cos\Bigl(-\frac{5\pi}{2}\Bigr) + i\sin\Bigl(-\frac{5\pi}{2}\Bigr)\Bigr) & $10 \cdot (-\frac{\pi}{4}) = -\frac{5\pi}{2}$. \\[3ex]
= 32\Bigl(\cos\Bigl(-\frac{\pi}{2}\Bigr) + i\sin\Bigl(-\frac{\pi}{2}\Bigr)\Bigr) = -32i & $(\sqrt2)^{10} = 2^5$, en $-\frac{5\pi}{2} + 2\pi = -\frac{\pi}{2}$.
\end{stappen}}

De formule werkt alleen voor de goniometrische vorm. In cartesische vorm mag je niet elk deel apart tot de macht verheffen: $(2 + 2i)^3$ is niet $8 + 8i$. Met De Moivre is $2 + 2i = 2\sqrt2\bigl(\cos\frac{\pi}{4} + i\sin\frac{\pi}{4}\bigr)$, dus
\[ (2 + 2i)^3 = \bigl(2\sqrt2\bigr)^3\Bigl(\cos\frac{3\pi}{4} + i\sin\frac{3\pi}{4}\Bigr) = 16\sqrt2\Bigl(-\frac{\sqrt2}{2} + \frac{\sqrt2}{2}i\Bigr) = -16 + 16i . \]
Ook de modulus mag je niet vergeten: $\bigl(2\sqrt2\bigr)^3 = 16\sqrt2$, en niet $2\sqrt2$.

\begin{werkwijze}
Schrijf het grondtal in goniometrische vorm. Neem de modulus tot de macht $n$ en vermenigvuldig het argument met $n$. Breng het argument terug naar $]-\pi, \pi]$ door een veelvoud van $2\pi$ op te tellen of af te trekken, en zet het resultaat in cartesische vorm.
\end{werkwijze}

\nujij{Bereken $(\sqrt3 + i)^3$.}{$\sqrt3 + i = 2\bigl(\cos\frac{\pi}{6} + i\sin\frac{\pi}{6}\bigr)$, dus $(\sqrt3 + i)^3 = 8\bigl(\cos\frac{\pi}{2} + i\sin\frac{\pi}{2}\bigr) = 8i$.}

\begin{oefeningen}
\oef{h4-moivre}

\oef{h4-goniomacht}
\end{oefeningen}

\subsection{Machten op een spiraal}

\naast{\begin{minipage}{8.8cm}
Teken je de opeenvolgende machten $1, z, z^2, z^3, \ldots$, dan draait elk punt over hetzelfde argument $\theta$ verder, en wordt de afstand tot de oorsprong telkens met $r$ vermenigvuldigd. Is $r > 1$, dan liggen de punten op een spiraal die naar buiten draait, zoals in de figuur voor $z = 1{,}15\bigl(\cos\frac{\pi}{6} + i\sin\frac{\pi}{6}\bigr)$. Is $r < 1$, dan draait de spiraal naar de oorsprong toe. Is $r = 1$, dan blijven alle machten op de cirkel met straal 1.

Bij $1 + i$ uit de Verken van hoofdstuk 3 was $r = \sqrt2$ en $\theta = \frac{\pi}{4}$: na acht stappen ben je een hele toer rond, op $(\sqrt2)^8 = 16$ keer de afstand.
\end{minipage}}{\begin{tikzpicture}[x=.58cm,y=.58cm,font=\scriptsize]
  \draw[->,grijs] (-4.2,0) -- (6,0) node[right,text=black] {Re};
  \draw[->,grijs] (0,-2.4) -- (0,3.6) node[above,text=black] {Im};
  \draw[accent!45,thick] plot[domain=0:12,samples=80,smooth] ({1.15^\x*cos(30*\x)},{1.15^\x*sin(30*\x)});
  \foreach \n in {0,...,12}{\fill[accent] ({1.15^\n*cos(30*\n)},{1.15^\n*sin(30*\n)}) circle (1.8pt);}
  \node[below right,inner sep=1pt] at (1,0) {$1$};
  \node[above right,inner sep=1pt] at ({1.15*cos(30)},{1.15*sin(30)}) {$z$};
  \node[above right,inner sep=1pt] at ({1.3225*cos(60)},{1.3225*sin(60)}) {$z^2$};
  \node[above,inner sep=2pt] at ({1.15^12},0) {$z^{12}$};
\end{tikzpicture}}

\voorbeeld{}{Bereken $z^2$, $z^3$ en $z^4$ voor $z = \frac12 + \frac12 i$, en teken $1, z, z^2, z^3, z^4$.\label{vb:h4-spiraal}}{%
\uitleg{De modulus van $z$ is kleiner dan 1: we verwachten een spiraal naar binnen. Met De Moivre vinden we elke macht meteen.}
\naast{\begin{minipage}{10.5cm}
\begin{stappenvrij}{6.2cm}{4cm}
z = \tfrac{\sqrt2}{2}\bigl(\cos\tfrac{\pi}{4} + i\sin\tfrac{\pi}{4}\bigr) & $r = \sqrt{\tfrac14 + \tfrac14} = \tfrac{\sqrt2}{2} < 1$. \\[1ex]
z^2 = \tfrac12\bigl(\cos\tfrac{\pi}{2} + i\sin\tfrac{\pi}{2}\bigr) = \tfrac12 i & $\bigl(\tfrac{\sqrt2}{2}\bigr)^2 = \tfrac12$. \\[1ex]
z^3 = \tfrac{\sqrt2}{4}\bigl(\cos\tfrac{3\pi}{4} + i\sin\tfrac{3\pi}{4}\bigr) = -\tfrac14 + \tfrac14 i & $\bigl(\tfrac{\sqrt2}{2}\bigr)^3 = \tfrac{\sqrt2}{4}$. \\[1ex]
z^4 = \tfrac14(\cos\pi + i\sin\pi) = -\tfrac14 &
\end{stappenvrij}

\vspace{1mm}
Bij elke stap draait het punt $\frac{\pi}{4}$ verder, en komt het $\frac{\sqrt2}{2} \approx 0{,}71$ keer zo dicht bij de oorsprong.
\end{minipage}}{\begin{tikzpicture}[x=2cm,y=2cm,font=\scriptsize]
  \draw[black!14,very thin,step=.25] (-.5,-.25) grid (1.25,.75);
  \draw[->,grijs] (-.5,0) -- (1.35,0) node[right,text=black] {Re};
  \draw[->,grijs] (0,-.25) -- (0,.85) node[above,text=black] {Im};
  \foreach \t in {.5,1}{\node[below,inner sep=1.5pt] at (\t,0) {$\t$};}
  \draw[accent!45,thick] plot[domain=0:4,samples=40,smooth] ({.7071^\x*cos(45*\x)},{.7071^\x*sin(45*\x)});
  \cpunt[above]{1}{0}{$1$}
  \cpunt[above right]{.5}{.5}{$z$}
  \cpunt[above right]{0}{.5}{$z^2$}
  \cpunt[above left]{-.25}{.25}{$z^3$}
  \cpunt[below left]{-.25}{0}{$z^4$}
\end{tikzpicture}}}

\begin{oefeningen}
\oef{h4-spiraal}
\end{oefeningen}

%========================================================================================
\section{Machtswortels}

\subsection{De \texorpdfstring{$n$}{n}-de machtswortels}

In de Verken vond je dat $\cos\frac{\pi}{6} + i\sin\frac{\pi}{6}$ een derdemachtswortel is van $i$, maar ook $-i$, want $(-i)^3 = -i^3 = i$. Dat komt doordat een argument maar vastligt op een veelvoud van $2\pi$ na: het argument van $i$ is $\frac{\pi}{2}$, maar ook $\frac{\pi}{2} + 2\pi$ en $\frac{\pi}{2} + 4\pi$. Delen door 3 geeft drie verschillende hoeken: $\frac{\pi}{6}$, $\frac{5\pi}{6}$ en $\frac{3\pi}{2}$.

\begin{theorie}
Een complex getal $w$ is een \textbf{\textcolor{accent}{$n$-de machtswortel}} van $z$ als \quad $\formule{w^n = z}$ \qquad ($n \geqslant 1$ natuurlijk)

\begin{stelling}{De $n$-de machtswortels}
Een complex getal $z = r(\cos\theta + i\sin\theta) \neq 0$ heeft precies $n$ verschillende $n$-de machtswortels:

\vspace{1mm}
\qquad $\formule{w_k = \sqrt[n]{r}\,\Bigl(\cos\frac{\theta + k \cdot 2\pi}{n} + i\sin\frac{\theta + k \cdot 2\pi}{n}\Bigr)}$ \qquad met $k = 0, 1, \ldots, n - 1$
\end{stelling}

\begin{bewijs}
Het getal 0 is geen machtswortel, want $0^n = 0 \neq z$. Zoek dus $w = s(\cos\varphi + i\sin\varphi)$ met $s > 0$. Volgens De Moivre is $w^n = s^n(\cos n\varphi + i\sin n\varphi)$. Twee getallen in goniometrische vorm zijn gelijk als hun moduli gelijk zijn en hun argumenten een veelvoud van $2\pi$ verschillen, dus $w^n = z$ als en slechts als
\bewijsregel{s^n = r \quad \text{en} \quad n\varphi = \theta + k \cdot 2\pi \text{ voor een geheel getal } k ,}
dus $s = \sqrt[n]{r}$ en $\varphi = \frac{\theta + k \cdot 2\pi}{n}$. De machtswortels zijn dus precies de getallen $w_k$, met $k$ geheel. Voor $k = 0, 1, \ldots, n - 1$ zijn ze verschillend: twee van die hoeken verschillen $(k - l)\frac{2\pi}{n}$ met $0 < |k - l| < n$, dus minder dan $2\pi$ en niet 0, en dat is geen veelvoud van $2\pi$. Elke andere $k$ geeft er een van: schrijf $k = qn + m$ met $m \in \{0, 1, \ldots, n - 1\}$, dan is $\frac{\theta + k \cdot 2\pi}{n} = \frac{\theta + m \cdot 2\pi}{n} + q \cdot 2\pi$, dus $w_k = w_m$.
\end{bewijs}
\end{theorie}

Alle $n$-de machtswortels hebben dezelfde modulus $\sqrt[n]{r}$, en hun argumenten liggen telkens $\frac{2\pi}{n}$ uit elkaar. Voor $n = 2$ vind je de twee vierkantswortels uit hoofdstuk 2 terug: $w$ en $-w$, een halve toer uit elkaar.

Let op: $\sqrt[n]{r}$ is hier de gewone reële wortel van het positieve getal $r$. Het wortelteken gebruiken we niet voor de complexe wortels: $8$ heeft drie derdemachtswortels, maar $\sqrt[3]{8}$ is alleen $2$. Wie de vergelijking $z^3 = -8$ oplost met alleen $z = -2$, mist er dus twee.

\voorbeeld{}{Los op in $\C$: $z^4 + 16 = 0$.\label{vb:h4-vgl}}{%
\uitleg{De oplossingen zijn de vierdemachtswortels van $-16$. Het argument van $-16$ is $\pi$, en de hoeken liggen telkens $\frac{2\pi}{4} = \frac{\pi}{2}$ uit elkaar.}
\begin{stappenbreed}
z^4 = -16 = 16(\cos\pi + i\sin\pi) & $r = 16$ en $\theta = \pi$. \\[1ex]
z_k = 2\Bigl(\cos\Bigl(\frac{\pi}{4} + k\frac{\pi}{2}\Bigr) + i\sin\Bigl(\frac{\pi}{4} + k\frac{\pi}{2}\Bigr)\Bigr) & $\sqrt[4]{16} = 2$, $k = 0, 1, 2, 3$. \\[3ex]
z = \sqrt2 + i\sqrt2 \quad \text{of} \quad z = -\sqrt2 + i\sqrt2 & De hoeken $\frac{\pi}{4}$ en $\frac{3\pi}{4}$, \\[1ex]
z = -\sqrt2 - i\sqrt2 \quad \text{of} \quad z = \sqrt2 - i\sqrt2 & en $\frac{5\pi}{4}$ en $\frac{7\pi}{4}$: een vierkant.
\end{stappenbreed}}

\voorbeeld{}{Bepaal de derdemachtswortels van $8i$, en stel ze voor in het vlak van Gauss.\label{vb:h4-wortel}}{%
\uitleg{We schrijven $8i$ in goniometrische vorm, nemen de reële derdemachtswortel van de modulus, en delen het argument plus $k \cdot 2\pi$ door 3.}
\begin{stappenbreed}
8i = 8\Bigl(\cos\frac{\pi}{2} + i\sin\frac{\pi}{2}\Bigr) & $r = 8$ en $\theta = \frac{\pi}{2}$. \\[3ex]
w_k = 2\Bigl(\cos\Bigl(\frac{\pi}{6} + k\frac{2\pi}{3}\Bigr) + i\sin\Bigl(\frac{\pi}{6} + k\frac{2\pi}{3}\Bigr)\Bigr) & $\sqrt[3]{8} = 2$, en $\frac{\pi}{2} : 3 = \frac{\pi}{6}$. \\[3ex]
w_0 = 2\Bigl(\cos\frac{\pi}{6} + i\sin\frac{\pi}{6}\Bigr) = \sqrt3 + i & $k = 0$. \\[3ex]
w_1 = 2\Bigl(\cos\frac{5\pi}{6} + i\sin\frac{5\pi}{6}\Bigr) = -\sqrt3 + i & $k = 1$: $\frac{\pi}{6} + \frac{4\pi}{6}$. \\[3ex]
w_2 = 2\Bigl(\cos\frac{3\pi}{2} + i\sin\frac{3\pi}{2}\Bigr) = -2i & $k = 2$: $\frac{\pi}{6} + \frac{8\pi}{6}$.
\end{stappenbreed}
\naast{\begin{minipage}{10cm}
Controle: $(-2i)^3 = -8i^3 = 8i$. De drie wortels liggen op de cirkel met straal 2, en vormen een gelijkzijdige driehoek. Het argument $\frac{3\pi}{2}$ van $w_2$ is geen hoofdwaarde, maar dat hoeft hier niet: het getal is hetzelfde.
\end{minipage}}{\begin{tikzpicture}[x=.7cm,y=.7cm,font=\small]
  \gaussleeg{-2}{2}{-2}{2}
  \draw[hulplijn] (0,0) circle (2);
  \draw[accent!60] (1.732,1) -- (-1.732,1) -- (0,-2) -- cycle;
  \cpunt[above right]{1.732}{1}{$w_0$}
  \cpunt[above left]{-1.732}{1}{$w_1$}
  \cpunt[below right]{0}{-2}{$w_2$}
\end{tikzpicture}}}

\begin{werkwijze}
Schrijf $z$ in goniometrische vorm. De modulus van elke wortel is $\sqrt[n]{r}$. Het eerste argument is $\frac{\theta}{n}$; tel er telkens $\frac{2\pi}{n}$ bij tot je $n$ wortels hebt. Zet ze in cartesische vorm als de hoeken bijzonder zijn, en controleer één wortel door ze tot de $n$-de macht te verheffen.
\end{werkwijze}

\nujij{Bepaal de vierkantswortels van $-9$ met deze methode.}{$-9 = 9(\cos\pi + i\sin\pi)$: modulus 3, argumenten $\frac{\pi}{2}$ en $\frac{3\pi}{2}$, dus $3i$ en $-3i$.}

\begin{oefeningen}
\oef{h4-wortel}

\oef{h4-vgl}

\oef{h4-controle}
\end{oefeningen}

\subsection{Een regelmatige veelhoek}

In de voorbeelden vormden de derdemachtswortels een gelijkzijdige driehoek, en de vierdemachtswortels een vierkant. Dat is altijd zo.

\begin{theorie}
\begin{stelling}{Ligging van de machtswortels}
Voor $n \geqslant 3$ zijn de $n$-de machtswortels van een complex getal $z \neq 0$ de hoekpunten van een \textbf{\textcolor{accent}{regelmatige $n$-hoek}} met middelpunt 0.
\end{stelling}

\begin{bewijs}
Alle $n$ wortels hebben modulus $\sqrt[n]{r}$: ze liggen op de cirkel met middelpunt 0 en straal $\sqrt[n]{r}$. Hun argumenten zijn $\frac{\theta}{n}$, $\frac{\theta}{n} + \frac{2\pi}{n}$, \ldots, $\frac{\theta}{n} + (n - 1)\frac{2\pi}{n}$, en na nog een stap $\frac{2\pi}{n}$ ben je een volle toer rond, terug bij $w_0$. De $n$ punten verdelen de cirkel dus in $n$ bogen met dezelfde middelpuntshoek $\frac{2\pi}{n}$. Verbind elk punt met de oorsprong. Zo ontstaan $n$ gelijkbenige driehoeken met twee zijden $\sqrt[n]{r}$ en tophoek $\frac{2\pi}{n}$, die dus congruent zijn (ZHZ). Hun basissen, de zijden van de veelhoek, zijn dus even lang. Elke hoek van de veelhoek bestaat uit twee basishoeken van zulke driehoeken, dus alle hoeken zijn even groot.
\end{bewijs}
\end{theorie}

Voor $n = 2$ is er geen veelhoek: de twee vierkantswortels liggen tegenover elkaar, op een middellijn van de cirkel. Omdat opeenvolgende wortels $\frac{2\pi}{n}$ verschillen in argument, vind je de volgende wortel door de vorige te draaien over $\frac{2\pi}{n}$ om de oorsprong. Ken je één hoekpunt, dan vind je zo de hele veelhoek.

\voorbeeld{}{Een vierkant heeft middelpunt 0 en een hoekpunt in $3 + i$. Bepaal de andere hoekpunten, en de vergelijking $z^4 = w$ waarvan de hoekpunten de oplossingen zijn.\label{vb:h4-veelhoek}}{%
\uitleg{De hoekpunten zijn vierdemachtswortels van eenzelfde getal $w$. Opeenvolgende hoekpunten verschillen $\frac{2\pi}{4} = \frac{\pi}{2}$ in argument: we draaien telkens een kwartslag, dus vermenigvuldigen met $i$.}
\naast{\begin{minipage}{10.4cm}
\begin{stappenvrij}{5.9cm}{4.2cm}
(3 + i) \cdot i = -1 + 3i & Een kwartslag verder. \\[1ex]
(-1 + 3i) \cdot i = -3 - i & Nog een kwartslag: het tegengestelde van $3 + i$. \\[1ex]
(-3 - i) \cdot i = 1 - 3i & Het vierde hoekpunt. \\[1ex]
(3 + i)^2 = 9 + 6i + i^2 = 8 + 6i & Elk hoekpunt is een oplossing van $z^4 = w$, dus $w = (3 + i)^4$. \\[1ex]
(8 + 6i)^2 = 64 + 96i - 36 = 28 + 96i & De vergelijking is $z^4 = 28 + 96i$.
\end{stappenvrij}
\end{minipage}}{\begin{tikzpicture}[x=.48cm,y=.48cm]
  \gauss{-3}{3}{-3}{3}
  \draw[hulplijn] (0,0) circle (3.162);
  \draw[accent!60] (3,1) -- (-1,3) -- (-3,-1) -- (1,-3) -- cycle;
  \cpunt[right]{3}{1}{$3 + i$}
  \cpunt[above left]{-1}{3}{}
  \cpunt[left]{-3}{-1}{}
  \cpunt[below right]{1}{-3}{}
\end{tikzpicture}}}

\begin{oefeningen}
\oef{h4-veelhoek}
\end{oefeningen}

\subsection{Eenheidswortels}

De $n$-de machtswortels van 1 krijgen een eigen naam. Ze zijn bijzonder eenvoudig, want $1 = \cos 0 + i\sin 0$ heeft modulus 1 en argument 0.

\begin{theorie}
De $n$-de machtswortels van 1 heten de \textbf{\textcolor{accent}{$n$-de eenheidswortels}}:

\vspace{1mm}
\qquad $\formule{\cos\frac{k \cdot 2\pi}{n} + i\sin\frac{k \cdot 2\pi}{n}}$ \qquad met $k = 0, 1, \ldots, n - 1$
\end{theorie}

\naast{\begin{minipage}{10.2cm}
Ze liggen op de cirkel met straal 1, en vormen een regelmatige $n$-hoek met een hoekpunt in 1. De figuur toont de zesde eenheidswortels. Omdat de vergelijking $z^n = 1$ reële coëfficiënten heeft, komen de niet-reële eenheidswortels in toegevoegde paren, zoals je in hoofdstuk 2 bewees: de figuur is symmetrisch ten opzichte van de reële as. Zo hoef je maar de helft uit te rekenen.
\end{minipage}}{\begin{tikzpicture}[x=1.05cm,y=1.05cm,font=\scriptsize]
  \draw[->,grijs] (-1.4,0) -- (1.5,0) node[right,text=black] {Re};
  \draw[->,grijs] (0,-1.3) -- (0,1.4) node[above,text=black] {Im};
  \draw[hulplijn] (0,0) circle (1);
  \draw[accent!60] (0:1) \foreach \h in {60,120,...,300}{-- (\h:1)} -- cycle;
  \foreach \h in {0,60,...,300}{\fill[accent] (\h:1) circle (2pt);}
  \node[below right,inner sep=1.5pt] at (0:1) {$1$};
  \node[above right,inner sep=1.5pt] at (60:1) {$\frac12 + \frac{\sqrt3}{2}i$};
  \node[below left,inner sep=1.5pt] at (180:1) {$-1$};
\end{tikzpicture}}

\voorbeeld{}{Bepaal de achtste eenheidswortels in cartesische vorm.\label{vb:h4-eenheid}}{%
\uitleg{De argumenten zijn $k \cdot \frac{2\pi}{8} = k \cdot \frac{\pi}{4}$. We rekenen de eerste helft uit, en vinden de rest met de toegevoegde of het tegengestelde.}
\naast{\begin{minipage}{10.6cm}
\begin{stappenvrij}{5.6cm}{4.5cm}
k = 0, 2, 4, 6\colon \ 1, \ i, \ -1, \ -i & Op de assen: hoeken $0$, $\frac{\pi}{2}$, $\pi$, $\frac{3\pi}{2}$. \\[1ex]
k = 1\colon \ \tfrac{\sqrt2}{2} + \tfrac{\sqrt2}{2}i & Hoek $\frac{\pi}{4}$. \\[1ex]
k = 3\colon \ -\tfrac{\sqrt2}{2} + \tfrac{\sqrt2}{2}i & Hoek $\frac{3\pi}{4}$. \\[1ex]
k = 5, 7\colon \ -\tfrac{\sqrt2}{2} - \tfrac{\sqrt2}{2}i, \ \tfrac{\sqrt2}{2} - \tfrac{\sqrt2}{2}i & De toegevoegden van $k = 3$ en $k = 1$.
\end{stappenvrij}

\vspace{1mm}
De acht punten vormen een regelmatige achthoek op de cirkel met straal 1.
\end{minipage}}{\begin{tikzpicture}[x=.9cm,y=.9cm,font=\scriptsize]
  \draw[->,grijs] (-1.4,0) -- (1.5,0) node[right,text=black] {Re};
  \draw[->,grijs] (0,-1.3) -- (0,1.4) node[above,text=black] {Im};
  \draw[hulplijn] (0,0) circle (1);
  \draw[accent!60] (0:1) \foreach \h in {45,90,...,315}{-- (\h:1)} -- cycle;
  \foreach \h in {0,45,...,315}{\fill[accent] (\h:1) circle (2pt);}
  \node[below right,inner sep=1.5pt] at (0:1) {$1$};
  \node[above right,inner sep=1.5pt] at (90:1) {$i$};
\end{tikzpicture}}}

\begin{oefeningen}
\oef{h4-eenheid}
\end{oefeningen}

%========================================================================================
\opdrachtenkop

\begin{oefeningen}
\oef{h4-o16}

\oef{h4-o3}

\oef{h4-o20}

\oef{h4-o17}

\oef{h4-o2}

\oef{h4-o19}

\oef{h4-o5}

\oef{h4-o1}

\oef{h4-o4}

\oef{h4-o9}

\oef{h4-o18}

\oef{h4-o12}

\oef{h4-o6}

\oef{h4-o22}

\oef{h4-o23}

\oef{h4-o7}

\oef{h4-o24}

\oef{h4-o10}

\oef{h4-o8}

\oef{h4-o21}

\oef{h4-o11}

\oef{h4-o14}

\oef{h4-o13}

\oef{h4-o25}

\oef{h4-o15}

\oef{h4-o26}
\end{oefeningen}

%----------------------------------------------------------------------------------------
% KOPPEL: verbanden met andere leerstof
%----------------------------------------------------------------------------------------
\needspace{20\baselineskip}
\opdrachtsectie{Koppel}

\lkr{De eerste oefening gebruikt het merkwaardig product $(a + b)^3$. De tweede is een recursieve rij, uit de cursus Rijen (6de jaar): het begrip volstaat, de leerstof van die cursus is niet nodig. Het leerplan noemt fractalen als mogelijke uitbreiding.}

\begin{oefeningen}
\oef{h4-koppel-1}

\oef{h4-koppel-2}
\end{oefeningen}

%----------------------------------------------------------------------------------------
\opdrachtsectie{Test jezelf}

\begin{oefeningen}
\oef{h4-t1}

\oef{h4-t2}

\oef{h4-t7}

\oef{h4-t3}

\oef{h4-t4}

\oef{h4-t5}

\oef{h4-t6}
\end{oefeningen}

\herhaalschema{%
  \fv{oef:h4-t1}{vb:h4-moivre}{oef:h4-moivre}%
  \fv{oef:h4-t2}{vb:h4-moivre}{oef:h4-goniomacht}%
  \fv{oef:h4-t7}{vb:h4-spiraal}{oef:h4-spiraal}%
  \fv{oef:h4-t3}{vb:h4-wortel}{oef:h4-wortel}%
  \fv{oef:h4-t4}{vb:h4-vgl}{oef:h4-vgl}%
  \fv{oef:h4-t5}{vb:h4-moivre2}{oef:h4-moivre}%
  \fv{oef:h4-t6}{vb:h4-veelhoek}{oef:h4-veelhoek}}

\ifzelfstudie
\section*{Maak je eigen samenvatting}
Sluit dit boek en schrijf in je schrift: de formule van De Moivre en hoe je ze bewijst, hoe je een hoge macht berekent, waarom opeenvolgende machten op een spiraal liggen, wat een $n$-de machtswortel is, de formule voor de $n$-de machtswortels en waarom er precies $n$ zijn, hoe ze in het vlak liggen, en wat eenheidswortels zijn. Vergelijk daarna met de kaders.
\fi
